\documentclass[11pt]{article}
\usepackage[utf8]{inputenc}
\usepackage[T1]{fontenc}
\usepackage{lmodern}
\usepackage{amsmath,amssymb,amsthm}
\usepackage{geometry}
\usepackage{hyperref}
\usepackage{booktabs}
\geometry{margin=1in}

\title{Turing Laboratory Execution Protocol\\
Version 1.0}
\author{M. Drake Stapleton\\AIEN Project}
\date{September 29, 2026}

\begin{document}

\maketitle

\noindent\textit{Engineering companion to ``Turing Instrument Calibration and Validation Program'' (v1.0). The calibration program defines the experiment order and scientific gates: EXP-001 validates the compression bridge; EXP-002A/B/C test false structure, additional complexity, and mechanism discrimination; EXP-002D tests unification; and EXP-003 tests active experiment selection. This document specifies the repository organization, schemas, command names, and statistical implementation that make those experiments executable.}

\section*{Scope}

This document converts the Turing calibration plan into an executable laboratory program for:
\begin{verbatim}
CAL-0
Measurement freeze

EXP-001
Actual compression bridge

EXP-002A
Pure diffusion / false-structure rejection

EXP-002B
Drift / Occam transition

EXP-002C
Mean-reversion / mechanism discrimination

EXP-002D
Unification

EXP-003
Active experiment selection

H3
Replay and statistical reproducibility
\end{verbatim}
H4 energy closure and H5 Turing yield should use the same artifact and receipt architecture but remain gated until the informational measurement has passed EXP-001 through EXP-003.

The laboratory must behave as though a skeptical outside group will receive the complete artifact bundle and attempt to reproduce every conclusion without access to the original implementation.

\part{Repository and Trust Structure}

\section{Canonical repository layout}

Recommended top-level repository:
\begin{verbatim}
turing-lab/
|
|-- Cargo.toml
|-- README.md
|-- LICENSE
|-- rust-toolchain.toml
|
|-- profiles/
|   |-- schema/
|   |   `-- turing-profile.schema.json
|   |
|   |-- draft/
|   |   `-- Turing-profile-v1.0.toml
|   |
|   `-- frozen/
|       |-- Turing-profile-v1.0.toml
|       |-- Turing-profile-v1.0.sha256
|       `-- FREEZE_RECEIPT.json
|
|-- crates/
|   |-- turing-core/
|   |-- turing-artifact/
|   |-- turing-profile/
|   |-- turing-model-code/
|   |-- turing-probability/
|   |-- turing-score/
|   |-- turing-score-independent/
|   |-- turing-coder-range/
|   |-- turing-coder-rans/
|   |-- turing-generator/
|   |-- turing-stats/
|   |-- turing-blind/
|   |-- turing-receipt/
|   |-- turing-verify/
|   `-- turing-cli/
|
|-- experiments/
|   |-- EXP-001/
|   |-- EXP-002A/
|   |-- EXP-002B/
|   |-- EXP-002C/
|   |-- EXP-002D/
|   `-- EXP-003/
|
|-- generators/
|   |-- diffusion/
|   |-- drift_diffusion/
|   |-- ornstein_uhlenbeck/
|   `-- experiment_selection/
|
|-- baselines/
|   |-- uniform/
|   |-- empirical/
|   |-- markov/
|   `-- fixed_heuristic/
|
|-- controls/
|   |-- memorizer/
|   |-- polynomial/
|   `-- flexible_predictor/
|
|-- scripts/
|   |-- power_simulation.py
|   |-- verify_no_overlap.py
|   |-- verify_candidate_freeze.py
|   |-- seal_manifest.py
|   `-- summarize_results.py
|
|-- schemas/
|   |-- candidate.schema.json
|   |-- dataset.schema.json
|   |-- probability-stream.schema.json
|   |-- code-result.schema.json
|   |-- score-result.schema.json
|   |-- experiment-receipt.schema.json
|   |-- intervention.schema.json
|   `-- replication.schema.json
|
|-- preregistration/
|   |-- EXP-001.md
|   |-- EXP-002A.md
|   |-- EXP-002B.md
|   |-- EXP-002C.md
|   |-- EXP-002D.md
|   `-- EXP-003.md
|
|-- sealed/
|   `-- .gitignore
|
|-- published/
|   `-- README.md
|
`-- docs/
    |-- MODEL_CODE.md
    |-- BLINDING.md
    |-- STATISTICS.md
    |-- CODING.md
    |-- REPLICATION.md
    |-- FAILURE_POLICY.md
    `-- CLAIM_LADDER.md
\end{verbatim}
The \texttt{sealed/} directory must never be available to candidate-development processes. In production, it should ideally live on a different account, process boundary, or physical machine.

\section{Process separation}

Define four security principals.
\begin{verbatim}
candidate
evaluator
verifier
replicator
\end{verbatim}
The candidate principal may read:
\begin{verbatim}
profile
development data
public generators where permitted
candidate APIs
\end{verbatim}
The candidate principal may not read:
\begin{verbatim}
sealed data
sealed seeds
sealed parameters
evaluator manifests
future interventions
counterfactual outcomes
\end{verbatim}
The evaluator principal owns hidden generation and scoring execution. The verifier receives frozen artifacts after evaluation. The replicator should be able to reproduce the experiment from published protocol without importing production scorer code.

\section{Git discipline}

Use immutable tags for every executed experiment. Example:
\begin{verbatim}
turing-profile-v1.0
exp001-preregistered-v1
exp001-candidate-freeze-v1
exp001-evaluation-v1
exp001-published-v1
\end{verbatim}
Never rewrite those tags. A failed run receives its own permanent tag. Example:
\begin{verbatim}
exp002b-evaluation-v1-FAIL
\end{verbatim}
A corrected run is:
\begin{verbatim}
exp002b-preregistered-v2
\end{verbatim}
not a force-push over v1.

\part{Canonical Artifacts}

\section{Digest rule}

Every durable artifact receives:
\begin{verbatim}
sha256
byte_length
media_type
schema_version
created_by
created_at
\end{verbatim}
Canonical serialization should be deterministic. For JSON artifacts:
\begin{verbatim}
UTF-8
normalized numeric representation
sorted map keys
no insignificant whitespace
newline termination defined by profile
\end{verbatim}
Prefer a canonical JSON encoding or canonical CBOR for signed objects. Do not hash pretty-printed documents whose serialization may vary.

\section{Measurement profile schema}

Conceptual schema:
\begin{verbatim}
{
  "profile_id": "Turing-profile-v1.0",
  "profile_version": 1,
  "profile_digest": "sha256:...",

  "model_code": {
    "format": "turing-model-code-v1",
    "parameter_quantization": "...",
    "shared_background_digest": "sha256:...",
    "dependency_accounting": "explicit"
  },

  "probability": {
    "representation": "...",
    "quantization_bits": 0,
    "normalization": "...",
    "zero_probability_policy": "..."
  },

  "coding": {
    "coder_a": "range-v1",
    "coder_b": "rans-v1",
    "header_accounting": true,
    "side_information_accounting": true
  },

  "statistics": {
    "confidence_level": 0.95,
    "resampling_unit": "trajectory",
    "method": "...",
    "stopping_rule": "fixed-n"
  },

  "blinding": {
    "candidate_freeze_required": true,
    "sealed_release_after_freeze": true
  }
}
\end{verbatim}
The literal values must come from the preregistered methodology rather than this illustrative schema.

\section{Candidate manifest}

Every candidate receives a manifest before test release.
\begin{verbatim}
{
  "candidate_id": "M_diffusion_001",
  "candidate_family": "diffusion",
  "artifact_digest": "sha256:...",
  "model_code_digest": "sha256:...",
  "model_bits": 812,
  "dependencies": [
    {
      "digest": "sha256:...",
      "charged_bits": 0,
      "reason": "declared shared background"
    }
  ],
  "frozen_at": "2026-...",
  "profile_digest": "sha256:..."
}
\end{verbatim}
No mutable dependency may be referenced by branch name, URL, \texttt{latest}, or local path.

\section{Dataset manifest}

\begin{verbatim}
{
  "dataset_id": "EXP002A-test-001",
  "experiment": "EXP-002A",
  "split": "sealed_test",

  "generator_digest": "sha256:...",
  "parameter_commitment": "sha256:...",
  "seed_commitment": "sha256:...",

  "trajectory_count": 0,
  "trajectory_length": 0,

  "payload_digest": "sha256:...",
  "profile_digest": "sha256:..."
}
\end{verbatim}
During blinding, the parameter and seed values remain private. Their commitments are public. After the predetermined release condition, disclose them so anyone can verify the commitment.

\section{Probability-stream format}

The probability stream is one of the most important artifacts. The model produces it once. Coders consume it without invoking the model.

Header:
\begin{verbatim}
{
  "format": "turing-probability-stream-v1",
  "candidate_digest": "sha256:...",
  "dataset_digest": "sha256:...",
  "profile_digest": "sha256:...",
  "alphabet_digest": "sha256:...",
  "record_count": 0
}
\end{verbatim}
Record:
\begin{verbatim}
{
  "index": 1281,
  "context_digest": "sha256:...",
  "observed_symbol": 17,

  "probabilities": [
    ["0", "0.0010"],
    ["1", "0.0061"]
  ],

  "normalization_error": "0.0"
}
\end{verbatim}
For a continuous process, define the predictive coding representation explicitly. Possible approaches include:
\begin{verbatim}
fixed quantization bins
registered residual quantizer
registered density discretization
\end{verbatim}
but exactly one method must be frozen in the profile before evaluation. Never discretize differently for different candidates.

\section{Coder result}

\begin{verbatim}
{
  "coder": "range-v1",
  "probability_stream_digest": "sha256:...",
  "source_dataset_digest": "sha256:...",

  "payload_bits": 0,
  "header_bits": 0,
  "side_information_bits": 0,

  "total_bits": 0,

  "encoded_digest": "sha256:...",
  "decoded_digest": "sha256:...",

  "lossless_roundtrip": true
}
\end{verbatim}
All headers and required side information must be counted.

\section{Score record}

\begin{verbatim}
{
  "candidate_id": "M1",
  "baseline_id": "B0",

  "L_baseline_model": 0,
  "L_candidate_model": 0,

  "L_baseline_data_ideal": 0,
  "L_candidate_data_ideal": 0,

  "L_baseline_data_coder_a": 0,
  "L_candidate_data_coder_a": 0,

  "T_ideal": 0,
  "T_coder_a": 0,

  "uncertainty": {
    "lower": 0,
    "upper": 0,
    "method": "..."
  }
}
\end{verbatim}

\section{Final experiment receipt}

\begin{verbatim}
{
  "experiment_id": "EXP-001",
  "run_id": "EXP001-EVAL-0001",

  "profile_digest": "sha256:...",
  "preregistration_digest": "sha256:...",

  "candidate_root": "sha256:...",
  "dataset_root": "sha256:...",
  "probability_root": "sha256:...",
  "coding_root": "sha256:...",
  "analysis_root": "sha256:...",

  "primary_endpoint": "...",
  "gate": "PASS",

  "signed_by": "...",
  "signature": "..."
}
\end{verbatim}

\part{Common Command-Line Interface}

\section{CLI design}

Use a single control-plane CLI:
\begin{verbatim}
turing
\end{verbatim}
Suggested top-level commands:
\begin{verbatim}
turing profile
turing prereg
turing power
turing data
turing candidate
turing predict
turing code
turing score
turing verify
turing stats
turing experiment
turing replicate
turing publish
\end{verbatim}

\section{Profile commands}

Draft validation:
\begin{verbatim}
turing profile validate \
  profiles/draft/Turing-profile-v1.0.toml
\end{verbatim}
Freeze:
\begin{verbatim}
turing profile freeze \
  profiles/draft/Turing-profile-v1.0.toml \
  --out profiles/frozen/
\end{verbatim}
Verify:
\begin{verbatim}
turing profile verify \
  profiles/frozen/Turing-profile-v1.0.toml
\end{verbatim}
After freeze, editing the profile requires a new version.

\section{Power simulation}

Example:
\begin{verbatim}
turing power run \
  --experiment EXP-002B \
  --config preregistration/EXP-002B.power.toml \
  --replicates 10000 \
  --out experiments/EXP-002B/power/
\end{verbatim}
Output must include:
\begin{verbatim}
assumed effect sizes
sample size
estimated interval width
estimated false-positive rate
estimated false-negative rate
simulation seed
simulation code digest
\end{verbatim}
The output becomes part of preregistration.

\section{Candidate freeze}

\begin{verbatim}
turing candidate freeze \
  --candidate build/candidates/M1 \
  --profile profiles/frozen/Turing-profile-v1.0.toml \
  --out experiments/EXP-001/frozen_candidates/M1/
\end{verbatim}
Then:
\begin{verbatim}
turing candidate verify-freeze \
  experiments/EXP-001/frozen_candidates/M1/
\end{verbatim}
The evaluator must refuse any candidate whose freeze receipt is newer than the sealed-data release timestamp.

\section{Hidden data generation}

Evaluator-only command:
\begin{verbatim}
turing data generate \
  --experiment EXP-002A \
  --split sealed_test \
  --profile profiles/frozen/Turing-profile-v1.0.toml \
  --secret-config /secure/EXP002A-secret.toml \
  --out /secure/sealed/EXP-002A/
\end{verbatim}
Publish only commitments:
\begin{verbatim}
turing data commit \
  /secure/sealed/EXP-002A/ \
  --publish experiments/EXP-002A/commitments/
\end{verbatim}

\section{Prediction}

After candidate freeze:
\begin{verbatim}
turing predict \
  --candidate experiments/EXP-001/frozen_candidates/M1/ \
  --dataset /secure/sealed/EXP-001/test.bin \
  --profile profiles/frozen/Turing-profile-v1.0.toml \
  --out /secure/results/EXP-001/M1.prob
\end{verbatim}
Freeze the probability stream immediately.

\section{Coding}

Range coder:
\begin{verbatim}
turing code range \
  --probabilities M1.prob \
  --observations test.bin \
  --out M1.range
\end{verbatim}
rANS:
\begin{verbatim}
turing code rans \
  --probabilities M1.prob \
  --observations test.bin \
  --out M1.rans
\end{verbatim}
Verify:
\begin{verbatim}
turing code verify M1.range
turing code verify M1.rans
\end{verbatim}

\section{Scoring}

Primary scorer:
\begin{verbatim}
turing score primary \
  --candidate M1 \
  --baseline B0 \
  --dataset test.bin \
  --profile Turing-profile-v1.0.toml \
  --out score-primary.json
\end{verbatim}
Independent scorer:
\begin{verbatim}
turing score independent \
  --candidate-artifacts published/... \
  --out score-independent.json
\end{verbatim}
The independent scorer must be a separate crate/package with no dependency on \texttt{turing-score}.

\section{Gate execution}

\begin{verbatim}
turing experiment gate EXP-001 \
  --prereg preregistration/EXP-001.md \
  --artifacts experiments/EXP-001/run-0001/
\end{verbatim}
Output:
\begin{verbatim}
PASS
FAIL
INCONCLUSIVE
\end{verbatim}
\texttt{INCONCLUSIVE} must exist. Not every non-pass is proof of a false hypothesis.

\part{Statistical Analysis}

\section{Experimental unit}

For stochastic-process experiments, the primary independent unit should normally be the independently generated trajectory or hidden world rather than each time step. Observations inside a trajectory are correlated. Do not calculate confidence intervals by pretending every time sample is an independent replicate. This follows the calibration program's instruction to use block resampling for correlated observations rather than treating them as independent.

\section{Primary uncertainty procedure}

Recommended default for the first calibration series:
\begin{verbatim}
trajectory-level paired bootstrap
\end{verbatim}
when comparing models evaluated on the same hidden trajectories. For candidates A and B, calculate per-world:
\begin{verbatim}
delta_T_i = T_A,i - T_B,i
\end{verbatim}
Resample entire worlds/trajectories with replacement. Estimate:
\begin{verbatim}
median or mean delta_T
95% interval
probability delta_T > 0 under bootstrap distribution
\end{verbatim}
The exact estimator must be frozen before evaluation.

\section{Sequential dependence}

If each trajectory contains time-series dependence, either:
\begin{verbatim}
treat the entire trajectory as one resampling unit
\end{verbatim}
or use a preregistered moving/stationary block bootstrap when within-trajectory inference is required. Do not choose the block size after inspecting candidate differences.

\section{EXP-001 statistics}

Primary quantities:
\begin{verbatim}
coder overhead:
L_coded - L_ideal

rank agreement:
ideal ordering vs coder ordering

T difference:
T_coded - T_ideal
\end{verbatim}
Recommended tests:
\begin{verbatim}
Spearman rank correlation
paired trajectory bootstrap for discrepancy
predefined equivalence interval for coder overhead
\end{verbatim}
Do not use only a significance test. EXP-001 is fundamentally an equivalence/calibration problem. Define an acceptable discrepancy envelope first.

\section{Memorizer control}

Primary question:
\begin{verbatim}
Does M_mem achieve positive generalized T after L(M_mem)?
\end{verbatim}
Test across independent sealed worlds. Preregister expected sign. If the memorizer wins, investigate the accounting before claiming that memorization is scientifically rewarded.

\section{EXP-002A statistics}

For each candidate compare:
\begin{verbatim}
held-out T
development-to-held-out degradation
false-structure selection rate
\end{verbatim}
Primary contrasts:
\begin{verbatim}
diffusion vs drifted diffusion
diffusion vs polynomial
diffusion vs OU
diffusion vs memorizer
\end{verbatim}
Preregister whether the experiment uses:
\begin{verbatim}
family-wise correction
false-discovery-rate control
or hierarchical testing
\end{verbatim}
Do not decide after seeing which comparison succeeds.

\section{EXP-002A pass criterion}

A recommended preregistered form is:
\begin{verbatim}
P(diffusion has greater held-out delta_T than
each designated false-structure control)
must exceed declared criterion,
and the uncertainty interval on the relevant
paired contrast must exclude the failure region.
\end{verbatim}
Insert actual numerical criteria only after power simulation. Do not copy arbitrary \texttt{p < .05} conventions without checking experimental power.

\section{EXP-002B statistics}

The main object is the Occam curve:
\begin{verbatim}
mu
->
delta_T(drift - diffusion)
\end{verbatim}
Fit or summarize this curve without assuming a convenient parametric shape unless preregistered. Estimate:
\begin{verbatim}
zero-effect behavior
transition region
strong-effect behavior
\end{verbatim}
The critical null-world test is:
\begin{verbatim}
mu = 0
\end{verbatim}
where the drift parameter should not systematically earn its cost.

\section{EXP-002B crossing point}

Define:
\begin{verbatim}
mu*
\end{verbatim}
as the effect magnitude at which expected incremental T crosses zero. Estimate uncertainty for \texttt{mu*}. Because the crossing may be broad, report an interval rather than false precision.

\section{EXP-002C statistics}

Primary contrast:
\begin{verbatim}
OU
vs
diffusion
vs
drift
\end{verbatim}
across intervention strata:
\begin{verbatim}
near equilibrium
moderate displacement
far displacement
\end{verbatim}
Use a paired design: all candidate models evaluate the same generated trajectories. Test for:
\begin{verbatim}
model x intervention-regime interaction
\end{verbatim}
because the central scientific claim is structural generalization under changed initial conditions.

\section{EXP-002C explanatory criterion}

A candidate earns stronger explanatory support only if:
\begin{verbatim}
candidate frozen before intervention

predictive advantage survives intervention

advantage behaves in the direction predicted by
the proposed mechanism
\end{verbatim}
Passive compression alone does not establish causality or mechanism.

\section{EXP-002D statistics}

Compute cumulative net description length as environments are added. For specialists:
\begin{verbatim}
L_total_specialists(k)
\end{verbatim}
For unified model:
\begin{verbatim}
L_total_unified(k)
\end{verbatim}
Then:
\begin{verbatim}
delta_T_unification(k)
\end{verbatim}
Plot or tabulate this as a function of \texttt{k}. The desired pattern may legitimately be:
\begin{verbatim}
specialists better at small k

unified theory better after amortization
\end{verbatim}
Include heterogeneous negative-control environments.

\section{EXP-003 statistics}

Primary outcome:
\begin{verbatim}
observations-to-discrimination
\end{verbatim}
Secondary outcomes:
\begin{verbatim}
experiment count
elapsed experimental time
eventual confidence
T gained
\end{verbatim}
Compare active strategy against each frozen passive baseline on matched hidden worlds. Recommended analysis:
\begin{verbatim}
paired bootstrap of cost-to-threshold

survival/time-to-threshold analysis
if some runs never reach threshold

failure fraction
\end{verbatim}
Do not delete runs that fail to discriminate. They are part of the result.

\part{EXP-001 Runbook}

\section{Phase 001-A: engineering qualification}

Before sealed science:
\begin{verbatim}
range coder unit tests
rANS unit tests
round-trip tests
probability normalization tests
canonical serialization tests
digest tests
model-cost tests
candidate-freeze tests
dataset-overlap tests
independent-scorer parity on synthetic fixtures
\end{verbatim}
Required development gate:
\begin{verbatim}
EXP001_INFRA_READY = PASS
\end{verbatim}

\section{Phase 001-B: synthetic known-answer tests}

Create toy distributions where codelength is analytically known. Examples:
\begin{verbatim}
uniform binary source

strongly biased Bernoulli source

deterministic source with registered zero-probability handling

small Markov source
\end{verbatim}
Verify:
\begin{verbatim}
ideal codelength
~=
actual coder length
\end{verbatim}
inside known implementation overhead. This is engineering validation, not scientific EXP-001 evidence.

\section{Phase 001-C: candidate freeze}

Freeze:
\begin{verbatim}
B0
B1
B2
B3
M_candidate
M_mem
\end{verbatim}
Generate:
\begin{verbatim}
candidate_manifest.json
artifact digest
model-code artifact
L(M) receipt
freeze timestamp
\end{verbatim}
After this point, candidate code is immutable.

\section{Phase 001-D: sealed evaluation}

Evaluator:
\begin{verbatim}
loads candidate
loads sealed observations
produces probability artifact
hashes probability artifact
runs both coders
verifies both decoders
calculates primary scores
\end{verbatim}
No candidate optimization occurs.

\section{Phase 001-E: independent verification}

Verifier receives:
\begin{verbatim}
profile
candidate manifest
model-code artifacts
probability streams
observations
coder artifacts
\end{verbatim}
Verifier does not receive the primary scorer source as a dependency. It independently reconstructs the result.

\section{EXP-001 final report}

Required sections:
\begin{verbatim}
Preregistration

Blinding

Candidate freeze

Model-cost accounting

Probability stream

Ideal codelength

Actual coder results

Memorizer result

Independent verification

Uncertainty

Failures/deviations

Gate decision
\end{verbatim}
Required summary table:
\begin{verbatim}
candidate
L(M)
L_ideal
L_range
L_rANS
T_ideal
T_range
T_rANS
uncertainty
gate
\end{verbatim}

\part{EXP-002 Runbook}

\section{Hidden-world generation architecture}

All stochastic generators should implement one common interface:
\begin{verbatim}
trait HiddenWorld {
    type Parameters;
    type State;
    type Observation;

    fn sample_parameters(secret_rng: &mut Rng) -> Self::Parameters;

    fn initialize(
        params: &Self::Parameters,
        secret_rng: &mut Rng
    ) -> Self::State;

    fn step(
        state: &mut Self::State,
        params: &Self::Parameters,
        secret_rng: &mut Rng
    ) -> Self::Observation;
}
\end{verbatim}
Candidate code must never receive \texttt{Parameters}. It receives only the observation stream and public metadata allowed by the profile.

\section{EXP-002A generation}

Private generator:
\begin{verbatim}
dX = sqrt(2D) dW
\end{verbatim}
Secret fields:
\begin{verbatim}
D
x0
seed
\end{verbatim}
Public fields may include:
\begin{verbatim}
dt
observation units
quantization definition
\end{verbatim}
if preregistered. Generate:
\begin{verbatim}
development worlds
calibration worlds
sealed worlds
replication worlds
\end{verbatim}
from non-overlapping seed domains.

\section{EXP-002B generation}

Private generator:
\begin{verbatim}
dX = mu dt + sqrt(2D) dW
\end{verbatim}
Secret:
\begin{verbatim}
mu
D
x0
seed
\end{verbatim}
Stratify hidden worlds across preregistered \texttt{mu} levels. Do not encode the stratum into filenames visible to candidates. Use opaque world IDs.

\section{EXP-002C generation}

Private generator:
\begin{verbatim}
dX = -theta X dt + sigma dW
\end{verbatim}
Secret:
\begin{verbatim}
theta
sigma
x0
seed
\end{verbatim}
For intervention evaluation, evaluator chooses starting state according to the preregistered intervention stratum. Candidate cannot update parameters after intervention identities are disclosed unless the protocol explicitly defines online updating. For the first calibration series, prefer frozen candidates.

\section{EXP-002D environment family}

Represent environment families in a sealed manifest:
\begin{verbatim}
{
  "family_id": "UFAM-001",
  "worlds": [
    "opaque-world-001",
    "opaque-world-002",
    "opaque-world-003"
  ]
}
\end{verbatim}
Do not expose which world belongs to which generator family before evaluation where that information would trivialize the task.

\part{EXP-003 Runbook}

\section{Experiment API}

EXP-003 needs an explicit intervention API. Conceptual request:
\begin{verbatim}
{
  "experiment_id": "INT-0041",

  "action": {
    "initial_displacement": "...",
    "observation_horizon": "...",
    "sampling_interval": "...",
    "replicates": 4
  },

  "predictions": [
    {
      "hypothesis": "H0",
      "distribution_digest": "sha256:..."
    },
    {
      "hypothesis": "H1",
      "distribution_digest": "sha256:..."
    }
  ],

  "expected_information_gain": "...",

  "committed_at": "...",
  "commitment_digest": "sha256:..."
}
\end{verbatim}
The evaluator must reject any intervention outside the legal action space.

\section{Commit-before-reveal protocol}

EXP-003 round:
\begin{verbatim}
candidate proposes intervention
->
candidate provides predicted distributions
->
candidate provides expected information gain
->
request artifact is hashed
->
evaluator verifies legality
->
experiment executes
->
result artifact is generated
->
result is revealed to candidate
\end{verbatim}
This prevents retrospective prediction.

\section{Passive-baseline replay}

The evaluator should precompute or execute:
\begin{verbatim}
fixed schedule
random legal policy
regular non-adaptive schedule
\end{verbatim}
over matched hidden worlds. Use the same discrimination rule for every strategy.

\section{Discrimination rule}

Freeze a criterion such as:
\begin{verbatim}
posterior odds
likelihood ratio
predefined model-evidence difference
or Turing-based decision boundary
\end{verbatim}
before evaluation. The choice must not depend on which method performs best.

\section{Counterfactual archive}

For each active choice, evaluator may later compute:
\begin{verbatim}
what would each other legal experiment
have revealed?
\end{verbatim}
Store this in:
\begin{verbatim}
counterfactual/
\end{verbatim}
but keep it hidden from the live candidate. Use it after the run to evaluate the quality of experiment selection.

\part{Agent Assignments}

\section{Agent 1: Metrology Lead}

Owns:
\begin{verbatim}
Turing-profile-v1.0
L(M)
shared-background definition
quantization rules
gate definitions
\end{verbatim}
Must not write candidate models. Primary deliverables:
\begin{verbatim}
profile
MODEL_CODE.md
STATISTICS.md
claim ladder
\end{verbatim}

\section{Agent 2: Blinding and Data Custodian}

Owns:
\begin{verbatim}
sealed-data generator
secret seeds
parameter commitments
candidate/test separation
dataset manifests
\end{verbatim}
Must not work on candidate ranking logic. Primary deliverables:
\begin{verbatim}
generator
sealed manifests
commitments
overlap audit
\end{verbatim}

\section{Agent 3: Range-Coder Lead}

Owns only:
\begin{verbatim}
turing-coder-range
\end{verbatim}
Must not inspect the rANS implementation. Deliverables:
\begin{verbatim}
range coder
decoder
round-trip tests
coder receipt
\end{verbatim}

\section{Agent 4: rANS Lead}

Owns only:
\begin{verbatim}
turing-coder-rans
\end{verbatim}
Must not reuse the range coder implementation. Independent code paths are intentional.

\section{Agent 5: Primary Scoring Lead}

Owns:
\begin{verbatim}
turing-score
\end{verbatim}
Implements:
\begin{verbatim}
L(M)
ideal codelength
coded T
uncertainty interface
\end{verbatim}
Does not own independent verification.

\section{Agent 6: Independent Verification Lead}

Owns:
\begin{verbatim}
turing-score-independent
turing-verify
\end{verbatim}
May read the specification. Must not import primary scorer libraries. Its job is to disagree when there is a bug.

\section{Agent 7: Statistics Lead}

Owns:
\begin{verbatim}
power simulation
bootstrap implementation
equivalence tests
multiple-comparison strategy
Occam transition analysis
EXP-003 cost-to-threshold analysis
\end{verbatim}
Must freeze analysis code before sealed outcome inspection where feasible.

\section{Agent 8: EXP-001 Integrator}

Owns no primitive scoring/coder implementation. Coordinates artifacts from Agents 1--7. Runs:
\begin{verbatim}
EXP-001 engineering qualification
candidate freeze
sealed run
verification
gate
report
\end{verbatim}

\section{Agent 9: EXP-002A Lead}

Owns:
\begin{verbatim}
pure-diffusion experiment
false-structure controls
development/sealed result assembly
\end{verbatim}
Uses common generator and scorer interfaces.

\section{Agent 10: EXP-002B Lead}

Owns:
\begin{verbatim}
drift sweep
effect-size design
Occam curve
mu* estimation
\end{verbatim}
Coordinates with Statistics Lead before preregistration.

\section{Agent 11: EXP-002C Lead}

Owns:
\begin{verbatim}
OU worlds
intervention regimes
mechanism-generalization analysis
\end{verbatim}

\section{Agent 12: EXP-002D Lead}

Owns:
\begin{verbatim}
specialist vs unified representation comparison
environment scaling
negative-control unification suite
\end{verbatim}

\section{Agent 13: EXP-003 Lead}

Owns:
\begin{verbatim}
legal intervention API
active experiment policy interface
passive baselines
decision threshold
counterfactual analysis
\end{verbatim}
Must not have access to the evaluator's hidden generator identity.

\section{Agent 14: Adversarial Reviewer}

Has read access to everything except unreleased sealed secrets. Attempts to break:
\begin{verbatim}
blinding
L(M)
metadata accounting
seed isolation
coder overhead accounting
statistical assumptions
candidate/test separation
intervention commitment
\end{verbatim}
Produces:
\begin{verbatim}
ADVERSARIAL_REVIEW.md
\end{verbatim}
before final gate signing.

\section{Agent 15: Reproduction Lead}

Starts from a clean checkout. Receives only what an external laboratory would receive. Must reproduce:
\begin{verbatim}
EXP-001
one EXP-002 family
\end{verbatim}
without assistance from original implementation authors once the run begins.

\part{Orchestration}

\section{Parallelization map}

The safe parallel work before CAL-0 freeze is:
\begin{verbatim}
Lane A
Profile + model-code specification

Lane B
Range coder

Lane C
rANS coder

Lane D
Primary scorer

Lane E
Independent scorer

Lane F
Generator framework

Lane G
Statistics/power tooling

Lane H
Artifact schemas + receipts
\end{verbatim}
These lanes have intentionally separate file ownership. After CAL-0 freeze, the metrology rules stop changing. Then:
\begin{verbatim}
EXP-001 executes serially

EXP-002A/B/C generator preparation
may proceed in parallel

but final evaluation occurs only after
the necessary preceding qualification gate
\end{verbatim}
EXP-003 infrastructure may be developed early, but scientific execution should wait until the measurement/discrimination stack has demonstrated credibility.

\section{Pull-request ownership rule}

Every PR declares:
\begin{verbatim}
owned paths
non-owned paths
experiment dependency
expected receipt
\end{verbatim}
Example:
\begin{verbatim}
Owned:
crates/turing-coder-rans/**

Do not modify:
crates/turing-coder-range/**
crates/turing-score/**
profiles/frozen/**
\end{verbatim}
No agent may edit a frozen profile.

\part{Full Command Sequence}

\section{Clean qualification build}

\begin{verbatim}
git clone <repo> turing-lab-clean
cd turing-lab-clean

cargo test --workspace --locked
\end{verbatim}
Then:
\begin{verbatim}
turing profile validate \
  profiles/draft/Turing-profile-v1.0.toml

turing power run \
  --all-preregistered

turing profile freeze \
  profiles/draft/Turing-profile-v1.0.toml
\end{verbatim}
Verify:
\begin{verbatim}
turing profile verify \
  profiles/frozen/Turing-profile-v1.0.toml
\end{verbatim}
Expected gate:
\begin{verbatim}
TURING_PROFILE_V1_FROZEN = PASS
\end{verbatim}

\section{EXP-001 commands}

Freeze candidates:
\begin{verbatim}
turing candidate freeze --all \
  --experiment EXP-001
\end{verbatim}
Verify chronology:
\begin{verbatim}
turing blind verify \
  --experiment EXP-001
\end{verbatim}
Generate/release sealed test to evaluator:
\begin{verbatim}
turing data evaluator-open \
  --experiment EXP-001
\end{verbatim}
Predict:
\begin{verbatim}
turing predict all \
  --experiment EXP-001
\end{verbatim}
Code:
\begin{verbatim}
turing code all \
  --experiment EXP-001 \
  --coders range,rans
\end{verbatim}
Score:
\begin{verbatim}
turing score primary \
  --experiment EXP-001
\end{verbatim}
Independent verification:
\begin{verbatim}
turing score independent \
  --experiment EXP-001
\end{verbatim}
Statistics:
\begin{verbatim}
turing stats run \
  --experiment EXP-001
\end{verbatim}
Gate:
\begin{verbatim}
turing experiment gate EXP-001
\end{verbatim}

\section{EXP-002A commands}

\begin{verbatim}
turing experiment prepare EXP-002A

turing candidate freeze --all \
  --experiment EXP-002A

turing data evaluator-open \
  --experiment EXP-002A

turing predict all \
  --experiment EXP-002A

turing code all \
  --experiment EXP-002A

turing score primary \
  --experiment EXP-002A

turing stats run \
  --experiment EXP-002A

turing experiment gate EXP-002A
\end{verbatim}

\section{EXP-002B commands}

\begin{verbatim}
turing power run \
  --experiment EXP-002B

turing prereg verify EXP-002B

turing candidate freeze --all \
  --experiment EXP-002B

turing experiment run EXP-002B

turing stats occam-curve \
  --experiment EXP-002B

turing experiment gate EXP-002B
\end{verbatim}

\section{EXP-002C commands}

Observational run:
\begin{verbatim}
turing experiment run EXP-002C \
  --phase observational
\end{verbatim}
Freeze again only if the preregistered protocol explicitly calls for an observational-fitting stage followed by candidate freeze. Then:
\begin{verbatim}
turing experiment run EXP-002C \
  --phase intervention \
  --candidate-mode frozen
\end{verbatim}
Analyze:
\begin{verbatim}
turing stats mechanism \
  --experiment EXP-002C
\end{verbatim}
Gate:
\begin{verbatim}
turing experiment gate EXP-002C
\end{verbatim}

\section{EXP-002D commands}

\begin{verbatim}
turing experiment run EXP-002D \
  --strategy specialists

turing experiment run EXP-002D \
  --strategy unified

turing stats unification \
  --experiment EXP-002D

turing experiment gate EXP-002D
\end{verbatim}

\section{EXP-003 commands}

Create initial ambiguous world:
\begin{verbatim}
turing experiment open EXP-003
\end{verbatim}
Candidate proposes:
\begin{verbatim}
turing experiment propose \
  --experiment EXP-003 \
  --candidate active-agent
\end{verbatim}
Commit:
\begin{verbatim}
turing experiment commit-intervention \
  --request intervention.json
\end{verbatim}
Evaluator executes:
\begin{verbatim}
turing experiment execute-intervention \
  --commitment <digest>
\end{verbatim}
Reveal:
\begin{verbatim}
turing experiment reveal-result \
  --round <n>
\end{verbatim}
Repeat until the fixed threshold or stopping rule is reached. Run baselines:
\begin{verbatim}
turing experiment baseline \
  EXP-003 \
  --policy fixed

turing experiment baseline \
  EXP-003 \
  --policy random

turing experiment baseline \
  EXP-003 \
  --policy regular
\end{verbatim}
Analyze:
\begin{verbatim}
turing stats active-vs-passive \
  --experiment EXP-003
\end{verbatim}
Counterfactual audit:
\begin{verbatim}
turing experiment counterfactual \
  --experiment EXP-003
\end{verbatim}
Gate:
\begin{verbatim}
turing experiment gate EXP-003
\end{verbatim}

\part{Hard Acceptance Gates}

\section{CAL-0}

\begin{verbatim}
profile complete
model-code accounting complete
power simulation complete
sample size fixed
analysis fixed
blinding verified
digest frozen
\end{verbatim}
Required:
\begin{verbatim}
TURING_PROFILE_V1_FROZEN = PASS
\end{verbatim}

\section{EXP-001}

\begin{verbatim}
both coders round-trip
actual coding tracks ideal coding
candidate result survives coder overhead
memorizer behaves as preregistered
independent scorer reproduces result
fresh-seed result preserves conclusion
\end{verbatim}
Required:
\begin{verbatim}
EXP_001_COMPRESSION_BRIDGE = PASS
\end{verbatim}
The calibration program treats this as the first priority because the current result has not yet produced a correspondingly smaller real encoded file.

\section{EXP-002A}

\begin{verbatim}
correct stochastic structure generalizes
invented drift rejected
invented mean reversion rejected
curve-fitting decoy fails held-out test
memorizer fails
\end{verbatim}
Required:
\begin{verbatim}
EXP_002A_RANDOMNESS_DISCRIMINATION = PASS
\end{verbatim}

\section{EXP-002B}

\begin{verbatim}
drift rejected at mu=0
ambiguous region remains appropriately uncertain
strong drift earns complexity
Occam curve reproduces
\end{verbatim}
Required:
\begin{verbatim}
EXP_002B_OCCAM_TRANSITION = PASS
\end{verbatim}

\section{EXP-002C}

\begin{verbatim}
mean reversion distinguished
wrong mechanisms rejected
frozen law survives changed starting state
intervention behavior agrees with mechanism
\end{verbatim}
Required:
\begin{verbatim}
EXP_002C_MECHANISM_DISCRIMINATION = PASS
\end{verbatim}
The calibration program makes this intervention step important because predictive compression alone does not justify calling a representation explanatory; stronger claims require invariance, counterfactual, causal, or mechanistic evidence.

\section{EXP-002D}

\begin{verbatim}
real reuse earns its abstraction cost
unified theory eventually amortizes when warranted
forced unification loses on negative controls
\end{verbatim}
Required:
\begin{verbatim}
EXP_002D_UNIFICATION = PASS
\end{verbatim}

\section{EXP-003}

\begin{verbatim}
active selection uses same discrimination threshold
predictions committed before outcomes
interventions remain within legal space
active policy reduces preregistered cost
result repeats on fresh hidden worlds
\end{verbatim}
Required:
\begin{verbatim}
EXP_003_ACTIVE_EXPERIMENT_SELECTION = PASS
\end{verbatim}
The calibration program defines the core objective as beating a preregistered passive schedule in observations or energy needed to distinguish competing explanations.

\part{Lab Notebook}

\section{Every run gets a notebook record}

Required fields:
\begin{verbatim}
run_id
operator
date/time
host identity
git commit
dirty-tree status
profile digest
experiment digest
candidate digests
dataset commitment
environment variables affecting behavior
hardware state
unexpected events
result
receipt digest
\end{verbatim}
No undocumented manual correction is permitted.

\section{Deviations}

A deviation record must contain:
\begin{verbatim}
what differed from preregistration
when discovered
whether outcome was already visible
scientific impact
decision:
    continue
    invalidate
    rerun under new preregistration
\end{verbatim}
A deviation discovered after results are visible must never be silently corrected in place.

\part{Publication Package}

\section{EXP-001 publication}

Release:
\begin{verbatim}
frozen profile
preregistration
candidate manifests
model-code artifacts
probability streams
both compressed streams
decoders
decoded verification
primary score
independent score
statistics
gate receipt
failure logs
\end{verbatim}
If secrecy is no longer needed, release sealed observations and seed preimages.

\section{EXP-002 publication}

Release:
\begin{verbatim}
all of EXP-001 reproducibility machinery
hidden-world commitments
eventual parameter/seed disclosure
candidate predictions
world-level results
bootstrap inputs
control results
rank-reversal analysis
\end{verbatim}

\section{EXP-003 publication}

Additionally release:
\begin{verbatim}
every intervention proposal
predicted distributions
information-gain estimates
commitment timestamps
observed results
passive baseline trajectories
counterfactual audit
cost-to-discrimination analysis
\end{verbatim}

\part{First 30-Day Lab Target}

The project should not try to execute everything at once. The first laboratory target should be:
\begin{verbatim}
WEEK 1
Freeze measurement/profile design
Finish model-code specification
Finish artifact schemas

WEEK 2
Finish two independent coders
Finish primary + independent scorer
Complete synthetic known-answer tests

WEEK 3
Freeze EXP-001 candidates
Run power/sample-size checks
Perform blinded EXP-001

WEEK 4
Independent reconstruction
Fresh-seed replication
Publish EXP-001 calibration bundle
\end{verbatim}
Do not treat these calendar labels as scientific deadlines if qualification is incomplete. The gate matters more than the date.

\part{First Implementation PRs}

Recommended PR sequence:
\begin{verbatim}
PR 1
Canonical artifact schemas

PR 2
Turing-profile loader + freezer

PR 3
Model-description accounting

PR 4
Canonical probability stream

PR 5
Range coder

PR 6
rANS coder

PR 7
Primary scorer

PR 8
Independent scorer

PR 9
Blinding + candidate freeze

PR 10
Power/statistics package

PR 11
EXP-001 orchestrator

PR 12
EXP-001 adversarial controls

PR 13
EXP-001 sealed evaluation artifacts

PR 14
EXP-001 independent replication

PR 15
EXP-002 generator framework

PR 16
EXP-002A

PR 17
EXP-002B

PR 18
EXP-002C

PR 19
EXP-002D

PR 20
EXP-003 intervention protocol
\end{verbatim}
Do not let PR 11 implement another coder or another scoring formula. Those primitives must remain independently reviewable.

\part{Orchestrator Prompt}

Use the following prompt for the lead engineering agent.
\begin{verbatim}
You are the lead orchestrator for the TURING calibration laboratory.

Your responsibility is not to make the Turing succeed.

Your responsibility is to execute the preregistered experiments faithfully enough that a failure is scientifically meaningful.

Start by reading the frozen Turing measurement profile, experimental preregistrations, and current repository implementation.

Do not change any frozen profile or preregistered gate after sealed outcomes become visible.

Decompose implementation into non-overlapping lanes:

1. artifact schemas;
2. model-description accounting;
3. range coder;
4. rANS coder;
5. primary scorer;
6. independent scorer;
7. blinding/sealed-data machinery;
8. statistical/power analysis;
9. experiment integration;
10. adversarial review.

Use separate agents for the two coders and two scorers.

No implementation agent may silently reuse the supposedly independent implementation.

For every lane require:
- exact owned paths;
- tests;
- immutable commit;
- evidence receipt;
- independent review.

Execute scientific experiments serially through the qualification ladder:

CAL-0
-> EXP-001
-> EXP-002A
-> EXP-002B
-> EXP-002C
-> EXP-002D
-> EXP-003.

Development of later infrastructure may occur in parallel, but no downstream scientific claim may bypass an upstream failed gate.

Before each experiment:
- verify clean repository state;
- verify exact profile digest;
- verify preregistration digest;
- verify candidate freeze timestamps;
- verify no candidate/sealed-data overlap;
- verify sample size was fixed before sealed outcomes;
- verify analysis code and primary endpoints.

After each experiment:
- preserve raw artifacts;
- preserve failures;
- run independent verification;
- generate uncertainty analysis;
- generate PASS / FAIL / INCONCLUSIVE receipt;
- never weaken the gate after seeing the outcome.

The final objective is not a large positive Turing score.

The objective is a calibration record in which the instrument demonstrably:
- corresponds to physical compression;
- rejects memorization;
- rejects false stochastic structure;
- rewards additional structure only when it earns its cost;
- distinguishes mechanisms;
- survives intervention;
- rewards justified unification;
- and helps select informative experiments.

Do not declare the Turing calibrated until the corresponding gates have actually passed.
\end{verbatim}

\part{Definition of Done}

The first major laboratory milestone is complete only when an outside reviewer can:
\begin{verbatim}
clone the repository

verify Turing-profile-v1.0

inspect candidate freeze receipts

inspect sealed-data commitments

recompute L(M)

recompute ideal predictive codelength

decode both compressed streams

verify exact observation recovery

recompute coded T

run the independent scorer

rerun the statistical analysis

inspect the memorizer control

inspect all deviations

verify the gate receipt
\end{verbatim}
without asking the original experiment team how the result was produced.

That is the standard. A successful scientific result that cannot be reconstructed from its artifacts is not a successful calibration result.

The central methodological point remains the same as in the calibration program: the strongest evidence comes from preregistered cases in which the metric can fail (rejecting false drift, deterministic fits to noise, and wrong mean-reversion models) rather than from maximizing \(T\). The final credibility step is independent implementation and replication, where the same scientific candidates should win and lose for the same reasons even if exact numerical Turings differ.

\end{document}
